\clearpage
\thispagestyle{plain}
\phantomsection
\addcontentsline{toc}{chapter}{ABSTRAK}

{\centering\fontsize{12}{14}\selectfont\bfseries ABSTRAK\par}

\vspace{0.2cm}

\noindent
\textbf{\MakeUppercase{\JudulPA}}

\vspace{0.2cm}

\noindent
\textbf{Nama\hspace{1.45cm}: \NamaMahasiswa}\\
\textbf{NRP\hspace{1.65cm}: \NRP}\\
\textbf{Pembimbing\hspace{0.2cm}: \PembimbingSatu}\\
\textbf{\phantom{Pembimbing\hspace{0.2cm}: }\PembimbingDua}

\vspace{0.2cm}

Penyakit kardiovaskular merupakan penyebab utama kematian di dunia, dan aritmia sering terjadi tanpa gejala sehingga monitoring elektrokardiogram (ECG) secara kontinu diperlukan untuk deteksi dini. Sistem monitoring yang ada umumnya berarsitektur \textit{cloud-centric} yang berlatensi tinggi dan bergantung pada koneksi internet, sedangkan model deteksi aritmia jarang diuji pada batasan mikrokontroler dan \textit{pipeline} data dari \textit{edge} ke \textit{cloud} belum tahan terhadap gangguan jaringan. Penelitian ini bertujuan membangun sistem monitoring detak jantung single-lead berbasis \textit{edge computing} yang menjalankan preprocessing sinyal ECG dan klasifikasi aritmia secara \textit{on-device} pada ESP32-S3 dengan pendekatan TinyML, serta mengirim hasilnya melalui \textit{pipeline} MQTT yang dirancang tahan terhadap putusnya koneksi. Filter Butterworth dan deteksi R-\textit{peak} Pan-Tompkins dijalankan di mikrokontroler, model CNN \textit{hybrid} dengan cabang morfologi dan cabang ritme dikuantisasi menjadi INT8, dan detak yang belum terkirim ditampung pada bufer lokal di PSRAM. Pada tahap progres ini, model \textit{float32} pada DS2 MIT-BIH dengan 22 pasien yang tidak terlihat saat pelatihan mencapai \textit{recall} 0,6996, F1 0,6594, dan AUC 0,9334. Kuantisasi INT8 menghasilkan model 22,94~KB dengan selisih \textit{recall} $+$0,0108, dan inferensinya di ESP32-S3 memakan 26,6~ms per detak. Dengan sinyal \textit{replay} dari rekaman MIT-BIH sebagai pengganti sinyal tubuh, \textit{pipeline} MQTT mencapai \textit{packet delivery success rate} 100,0\% dan 99,2\% pada dua ulangan dengan median latensi \textit{end-to-end} 939 dan 1.108~ms. Tujuan pertama tercapai, sedangkan dua tujuan lainnya berstatus Sebagian Tercapai karena uji putus koneksi pada konfigurasi akhir dan pengukuran daya belum selesai.
% sumber: bab4 tab:hasil-ds2 (DS2 22 pasien, recall 0,6996, F1 0,6594, AUC 0,9334), tab:delta-int8 (22,94 KB, selisih recall +0,0108), tab:sumber-daya (26,6 ms per detak), tab:target-pipeline (PDSR 100,0% dan 99,2%, median 939 dan 1.108 ms, dua ulangan, sumber sinyal replay), tab:ketercapaian (ketiga tujuan Sebagian Tercapai)

\vspace{0.25cm}
\noindent
\textbf{Kata kunci:} ECG, IoT, Edge Computing, TinyML, ESP32-S3, Aritmia, MQTT

\clearpage
